% 图契约：本题的源、边界和推导逐项对应；箭头符号按图注解释。
\documentclass[tikz,border=7pt]{standalone}
\usepackage{fontspec,xeCJK,amsmath,amssymb}
\IfFontExistsTF{Times New Roman}{\setmainfont{Times New Roman}}{\setmainfont{TeX Gyre Termes}}
\IfFontExistsTF{Heiti SC}{\setCJKmainfont[AutoFakeBold=1.5]{Heiti SC}}{\setCJKmainfont[AutoFakeBold=1.5]{FandolHei-Regular}}
\usetikzlibrary{arrows.meta,bending,calc,patterns,positioning}
\definecolor{ink}{HTML}{243745}
\definecolor{blue}{HTML}{275DAB}
\definecolor{green}{HTML}{177C63}
\definecolor{red}{HTML}{BD4935}
\definecolor{orange}{HTML}{AA6C16}
\definecolor{purple}{HTML}{715393}
\tikzset{every node/.style={font=\fontsize{9.3}{12.5}\selectfont,text=ink},
axis/.style={-{Stealth[length=4pt]},line width=.65pt,draw=ink!65},
vec/.style={-{Stealth[length=5pt]},line width=1pt,draw=blue},
short/.style={-{Stealth[length=3pt]},line width=.8pt},
guide/.style={dashed,draw=ink!40,line width=.5pt},
card/.style={rounded corners=2pt,draw=blue!65,fill=blue!4,align=center,text width=6.5cm,inner xsep=9pt,inner ysep=7pt},
title/.style={font=\fontsize{10.5}{14}\selectfont\bfseries,text width=15cm,align=center},
label/.style={font=\fontsize{8.4}{11}\selectfont,align=center},
note/.style={label,text width=6.2cm}}
\begin{document}
\begin{tikzpicture}[x=1cm,y=1cm]
\path[use as bounding box] (-8,-3.5) rectangle(8,3.5);
\node[title] at (0,3.15) {M3-028a\quad 初态与末态的电容和能量};
\begin{scope}[shift={(-4,0)}]

\foreach \x/\lab in {-1.3/{x=0},1.3/{x=L}}{
\begin{scope}[shift={(\x,0)}]
\draw[ink,thick] (-.8,-1)--(.8,-1) (-.8,1)--(.8,1);
\node[label] at(0,-1.45){$\lab$};\end{scope}}
\fill[blue!9] (.5,-.95) rectangle(2.1,.95);
\node at(-1.3,0){真空};\node at(1.3,0){$\varepsilon_r$};
\draw[vec,red] (-.8,1.8)--(.8,1.8) node[midway,above] {缓慢插入，$Q$ 固定};
\node[label] at(0,-2.2){两个端点的极板面积和间距完全相同};
\end{scope}
\node[card,draw=blue!65,fill=blue!4] at (3.9,1.8) {$C_L=\varepsilon_r C_0$};
\node[card,draw=green!65,fill=green!4] at (3.9,0) {$U_L=U_0/\varepsilon_r$};
\node[card,draw=orange!65,fill=orange!4] at (3.9,-1.8) {$W_{\rm field}=U_0(1-1/\varepsilon_r)$};
\end{tikzpicture}\end{document}